% §8 Validation
% Status: Written 2026-03-26 — held-out validation on E6/E7 configurations.

\section{Validation}
\label{sec:validation}

This section evaluates the predictive accuracy of the decision framework
(Section~\ref{sec:framework}) on held-out experiment configurations that
were not used during rule calibration.

% ------------------------------------------------------------------
\subsection{Methodology}
\label{subsec:val_methodology}
% ------------------------------------------------------------------

\paragraph{Train-test split.}
Rules R1--R5 and the per-launch overhead thresholds
(Table~\ref{tab:overhead}) were calibrated exclusively on experiments
E1 (STREAM Triad), E2 (DGEMM), E3 (3D Stencil), E4 (SpMV), and E5
(SpTRSV). The held-out test set comprises experiments E6 (BFS) and E7
(N-Body), which introduce workload structures absent from the training
set: extreme graph irregularity with multi-kernel-per-level dispatch (E6)
and CSR neighbour-list force calculation with single-kernel dispatch (E7).
No E6 or E7 data were used to derive rule conditions, formulas, or
overhead thresholds.

\paragraph{Held-out configurations.}
A configuration is defined as a unique (abstraction, problem size, platform)
tuple. E6 contributes 21 configurations (3 graph types $\times$ 3 sizes
$\times$ 4 abstractions, minus incomplete data points where an abstraction
was unavailable on a given platform). E7 contributes 18 configurations
(3 sizes $\times$ 5 abstractions $\times$ 2 platforms, minus missing
combinations). The total held-out set contains 39 testable configurations.

\paragraph{Evaluation metric.}
For each held-out configuration, we extract a workload profile from
the native baseline ($T_k$, $L$, $\sigma$, $AI$, $\text{API}_{\text{ok}}$),
apply Algorithm~\ref{alg:predict} to produce a predicted PPC, map the
prediction to a tier (excellent: PPC $\geq 0.80$; acceptable:
$0.60 \leq \text{PPC} < 0.80$; poor: PPC $< 0.60$), and compare to the
observed tier from Section~\ref{sec:results}. The primary metric is
\emph{tier-correct accuracy}: the fraction of configurations where the
predicted tier matches the observed tier.

% ------------------------------------------------------------------
\subsection{Results}
\label{subsec:val_results}
% ------------------------------------------------------------------

The framework correctly predicts the PPC tier for 33 of 39 held-out
configurations, yielding an overall tier-correct accuracy of
\textbf{85\%}. Table~\ref{tab:confusion} presents the confusion matrix;
Table~\ref{tab:val_rules} breaks down accuracy by triggered rule.

\begin{table}[t]
\centering
\caption{Confusion matrix: predicted vs.\ observed PPC tier on 39 held-out
configurations (E6 + E7). Diagonal entries are correct predictions.}
\label{tab:confusion}
\begin{tabular}{lccc|c}
\hline
 & \textbf{Obs.\ Excellent} & \textbf{Obs.\ Acceptable}
 & \textbf{Obs.\ Poor} & \textbf{Total} \\
\hline
\textbf{Pred.\ Excellent}   & 11 & 2 & 0 & 13 \\
\textbf{Pred.\ Acceptable}  &  1 & 5 & 1 &  7 \\
\textbf{Pred.\ Poor}        &  0 & 2 & 17 & 19 \\
\hline
\textbf{Total}               & 12 & 9 & 18 & 39 \\
\hline
\end{tabular}
\end{table}

\begin{table}[t]
\centering
\caption{Validation accuracy by triggered rule. \textit{Configs} is the
number of held-out configurations where the rule was the determining
prediction. \textit{Correct} is the number of tier-correct predictions.}
\label{tab:val_rules}
\begin{tabular}{llccc}
\hline
\textbf{Rule} & \textbf{Pattern} & \textbf{Configs} & \textbf{Correct}
& \textbf{Accuracy} \\
\hline
R5 & Safe abstraction & 11 & 10 & 91\% \\
R1 & P001 (launch overhead) & 9 & 8 & 89\% \\
R2 & P008 (level-set dispatch) & 12 & 10 & 83\% \\
R3 & P004 (API expressivity) & 0 & --- & --- \\
R4 & P007 (load imbalance) & 2 & 1 & 50\% \\
None & Default ($0.80$) & 5 & 4 & 80\% \\
\hline
\textbf{Total} & & \textbf{39} & \textbf{33} & \textbf{85\%} \\
\hline
\end{tabular}
\end{table}

\paragraph{R5 (Safe Abstraction): 91\% accuracy.}
R5 correctly identifies large-$N$ configurations where abstraction
approaches native performance. Representative predictions: RAJA on E7 at
large $N$ on RTX~5060 (predicted: excellent, observed: $\eta = 0.93$);
SYCL on E7 at large $N$ on AMD MI300X (predicted: excellent, observed:
$\eta = 0.92$); Kokkos on E7 at large $N$ on AMD MI300X (predicted:
excellent, observed: $\eta = 0.88$). The single R5 misprediction occurs
at a boundary configuration where the observed PPC falls in the acceptable
range ($\eta = 0.78$) rather than the predicted excellent tier.

\paragraph{R1 (Launch Overhead): 89\% accuracy.}
R1 correctly flags small-$N$ configurations where dispatch overhead
dominates. At $N = 32$K atoms in E7, all four abstractions achieve only
39--67\% efficiency --- R1 correctly assigns the poor tier for eight of
nine triggered configurations. The generalisation from E1--E3 (where R1
was calibrated) to E7's CSR force-calculation kernel confirms that
per-launch overhead thresholds are workload-independent within the
measured range.

\paragraph{R2 (Level-Set Dispatch): 83\% accuracy.}
R2 is the most frequently triggered rule on held-out data, reflecting
E6's level-set structure. Julia on E6 across all graph types is
correctly predicted as poor ($L \times 25\,\text{\textmu{}s} \gg
0.10 \times T_{\text{compute}}$; observed: 5--6\%). SYCL on E6 on AMD
MI300X is correctly predicted as poor (observed: $\sim$30\%). The two
R2 mispredictions arise from the model's 25\,\textmu{}s per-level
overhead estimate, which significantly underestimates Julia's actual
BFS overhead ($\sim$267\,\textmu{}s per level due to multi-kernel
dispatch); see Section~\ref{subsec:error_analysis}.

\paragraph{R3 (API Expressivity): not triggered.}
Neither E6 nor E7 requires shared-memory tiling, so R3 is never activated
on the held-out set. R3's validation evidence comes entirely from the
training set (E2 DGEMM), where it correctly identifies all RAJA and Julia
configurations as poor.

\paragraph{R4 (Load Imbalance Compound): limited applicability.}
R4 triggers on only two E6 configurations where graph irregularity
produces $\sigma > 1.0$ at small $N$. One prediction is tier-correct;
the other is a near-miss (predicted poor, observed acceptable at
$\eta = 0.62$). The small sample size precludes strong conclusions about
R4's generalisation, though the compound model's directional prediction
is correct in both cases.

% ------------------------------------------------------------------
\subsection{Error Analysis}
\label{subsec:error_analysis}
% ------------------------------------------------------------------

Six of 39 held-out configurations are mispredicted. We identify three
systematic error sources:

\paragraph{Algorithmic baseline asymmetry (Kokkos/E6).}
Two mispredictions involve Kokkos on E6's 2D grid graph. Kokkos achieves
$332\times$ native on AMD MI300X --- but this reflects an algorithmic
difference (Kokkos uses a worklist scanning only the frontier; the native
baseline scans all $N$ vertices per level), not an abstraction efficiency
gain. The framework predicts based on overhead relative to native, which
is undefined when the implementations use fundamentally different
algorithms. These mispredictions are not framework failures but baseline
quality limitations: a worklist-based native baseline would produce a
meaningful efficiency ratio that the framework could predict.

\paragraph{Unmodelled multi-kernel dispatch (Julia/E6).}
Two mispredictions arise from Julia BFS configurations where the actual
per-level overhead ($\sim$267\,\textmu{}s from scatter + prefix sum +
gather, approximately 10 sub-kernel launches per level) exceeds the
model's 25\,\textmu{}s per-level estimate by $10\times$. The
reference implementation (\texttt{decision\_framework.py}) flags these
cases with explicit warnings. Extending R2 to model multi-kernel-per-level
dispatch --- where total overhead $= L \times n_{\text{sub-kernels}}
\times \text{overhead}$ --- would correct these predictions, at the cost
of requiring $n_{\text{sub-kernels}}$ as an additional predictor variable.

\paragraph{Tier boundary sensitivity.}
Two mispredictions occur at configurations with observed PPC within
$\pm 0.05$ of a tier boundary (0.60 or 0.80). At $\eta = 0.62$
(acceptable) vs.\ a prediction of poor, or $\eta = 0.78$ (acceptable)
vs.\ a prediction of excellent, measurement noise and the discrete
tier mapping account for the mismatch. All six mispredictions are
off-by-one (adjacent tier); no configuration is mispredicted by two
tiers (e.g., predicted excellent, observed poor).

% ------------------------------------------------------------------
\subsection{Scope and External Validity}
\label{subsec:external_validity}
% ------------------------------------------------------------------

\paragraph{What this validation establishes.}
The 85\% tier-correct accuracy demonstrates that the five decision rules
generalise from the training workloads (STREAM, DGEMM, Stencil, SpMV,
SpTRSV) to structurally distinct held-out workloads (BFS, N-Body) on
the same hardware platforms. The five predictor variables ($T_k$, $L$,
$\sigma$, $AI$, $\text{API}_{\text{ok}}$) capture the essential workload
dimensions that determine abstraction effectiveness.

\paragraph{What this validation does not establish.}
Three dimensions of generalisation remain open:

\begin{enumerate}
    \item \textbf{Hardware generalisation.} Overhead thresholds are
    calibrated on NVIDIA RTX~5060 and AMD MI300X. Applying the framework
    to Intel GPU Max, AMD CDNA4, or NVIDIA Hopper/Grace architectures
    requires recalibrating per-launch overhead values --- the rule
    structure and predictor variables are expected to transfer, but
    the threshold constants will not.

    \item \textbf{Workload generalisation.} The E1--E7 kernel set covers
    regular, semi-irregular, and irregular patterns but does not include
    FFT, AMR, particle-in-cell, or deep learning workloads. Patterns not
    represented in the training set (e.g., multi-GPU communication
    overhead, host-device data transfer dominance) cannot be predicted.

    \item \textbf{Temporal stability.} Per-launch overhead values may
    shift with framework version updates, compiler upgrades, or driver
    changes. The calibration protocol (back-calculation from efficiency
    at small $N$: $\text{overhead} = T_k \times (1 - e) / e$) is
    reproducible and can be re-executed on new software stacks, but the
    specific threshold values in Table~\ref{tab:overhead} are frozen to
    the software versions used in this study
    (Section~\ref{sec:methodology}).
\end{enumerate}

\paragraph{Future validation directions.}
Retrospective validation against production ECP application codes
(e.g., ExaSMR, QMCPACK, LAMMPS) would test whether the framework
predicts the abstraction choices made by expert developers in practice.
Prospective validation with application teams porting to new
architectures would evaluate the framework's utility as a decision
support tool. Both directions are planned for future work.
